\subsection{(c) Normative OC Petri net $N_1$}
\label{sec:c}

\begin{figure}[H]
\centering
% Replace the placeholder below with: \includegraphics[width=\textwidth]{figures/N1_normative.pdf}
\placeholder[\textwidth]{%
  \textbf{N1\_normative.pdf}\\[4pt]
  Hand-drawn in draw.io / OC$\pi$.\\
  Export to PDF and place at \texttt{report/figures/N1\_normative.pdf}.
}
\caption{Normative OC Petri net $N_1$ for the Olympic transport process.
Transitions correspond to the seven event types; double (variable-arc) edges
on the \otPassenger{} place feeding \etDepart{} model the convergent batch
synchronization required by R(viii).}
\label{fig:n1}
\end{figure}

The net $N_1$ contains seven transitions, one for each event type in
\cref{tab:event-types}, connected by per-object-type places that carry tokens
representing live object instances.
The \otPassenger{} lifecycle begins at \etBooking, where a variable arc creates
one token per passenger enrolled; those tokens flow through \etIssue{} and
\etAssign{} before reaching an XOR split that routes each token to either
\etBoard{} or \etValidate{} depending on the \attrFare{} value (R(vii)).
Both branches reconverge at a shared \otPassenger{} place that feeds
\etDepart{} via a variable (double) arc, expressing the requirement that
\emph{all} passengers assigned to a departure must be present before the
transition fires (R(viii-a) convergence).
The \otTicket{} subnet runs in parallel: a ticket token is created at
\etIssue{} and consumed at the boarding transition, while \otBooking{} and
\otDeparture{} tokens are created by \etBooking{} and \etAssign{} respectively
and consumed when their lifecycle ends.
After \etDepart{} fires, individual \otPassenger{} tokens proceed independently
to \etComplete, reflecting the per-passenger nature of journey completion.
